A eq. (5.27) é derivada diretamente da Lei das Malhas de Kirchhoff (LMK) aplicada ao circuito da Figura 5.7.

\textbf{Lei das Malhas de Kirchhoff (LMK):}
A soma algébrica das quedas de tensão ao longo de qualquer malha fechada é igual à soma das tensões de fonte:
\[
\sum_{\rm malha} V = 0 \quad\text{(or}\quad V_{\rm fonte} = \sum V_{\rm elementos}).
\]

\textbf{Circuito RC da Figura 5.7:}
O circuito consiste em uma fonte $V_s$, um resistor $R$ e um capacitor $C$ em série. As quedas de tensão são:
\begin{itemize}
\item Resistor: $V_R = R\,i(t) = R\,C\,\dfrac{dV_C}{dt}$ (pois $i = C\,dV_C/dt$);
\item Capacitor: $V_C$.
\end{itemize}

\textbf{Aplicação da LMK:}
Percorrendo a malha no sentido da corrente:
\[
V_s = V_R + V_C = RC\,\frac{dV_C}{dt} + V_C.
\]
Reescrevendo:
\[
RC\,\frac{dV_C}{dt} + V_C = V_s.\quad\checkmark
\]
Esta é exatamente a eq. (5.27), confirmando que a Lei de Kirchhoff implica diretamente a equação diferencial do circuito RC. 